%% LyX 2.4.4 created this file.  For more info, see https://www.lyx.org/.
%% Do not edit unless you really know what you are doing.
\documentclass[english]{article}
\usepackage[T1]{fontenc}
\usepackage[latin9]{inputenc}
\usepackage{units}
\usepackage{enumitem}
\usepackage{amsmath}
\usepackage{amsthm}
\usepackage{geometry}
\geometry{verbose,tmargin=1in,bmargin=1in,lmargin=1in,rmargin=1in}
\PassOptionsToPackage{normalem}{ulem}
\usepackage{ulem}

\makeatletter
%%%%%%%%%%%%%%%%%%%%%%%%%%%%%% Textclass specific LaTeX commands.
\numberwithin{equation}{section}
\numberwithin{figure}{section}
\newlength{\lyxlabelwidth}      % auxiliary length 

%%%%%%%%%%%%%%%%%%%%%%%%%%%%%% User specified LaTeX commands.
\usepackage{fancyhdr}
\usepackage{lastpage}
\pagestyle{fancy}
\usepackage{enumitem}
\usepackage{ifthen}
\everymath=\expandafter{\the\everymath\displaystyle}
%\DeclareMathSizes{10}{10}{10}{10}
\usepackage{multicol}

\usepackage{tikz}
\usetikzlibrary{patterns}
\usetikzlibrary{plotmarks}
\usepackage{pgfplots}
\pgfplotsset{compat=newest}

\usepackage{advdate}    % Advancing/saving dates
\usepackage{datetime}   % Dates formatting
\usepackage{datenumber} % Counters for dates
\usdate
% Added by lyx2lyx
\setlength{\parskip}{\smallskipamount}
\setlength{\parindent}{0pt}

\makeatother

\usepackage{babel}
\begin{document}
\renewcommand{\author}{Dr.\,James Ashe}

\newcommand{\classNum}{232}

\newcommand{\classSec}{A and B}

\newcommand{\examType}{EXAM 1}

\renewcommand{\date}{\today}

%%First page's header%%%%%%%%%%%%%%%%%%%%%%
\fancypagestyle{firststyle} %%header settings for first pagr
{
	\fancyhead[L]{MTH \classNum{}(\classSec{}) \author{}\\
	\textbf{\examType{}} \date{}}
	\fancyhead[C]{\textbf{Name:}}
	\setlength{\headsep}{14pt}
}
%%%%%%%%%%%%%%%%%%%%%%%%%%%%%%%%%%%%%%%%%%

%%Next page's header%%%%%%%%%%%%%%%%%%%%%%%
\setlist{nolistsep}
\lhead{\classNum{}(\classSec{})} 
\chead{\textbf{Name:}} 
\cfoot{\thepage{}~of~\pageref{LastPage}} 
\setlength{\headsep}{5pt} 
%%%%%%%%%%%%%%%%%%%%%%%%%%%%%%%%%%%%%%%%%%%

%%Various settings; see the preamble for other settings%%%%%
%\setenumerate[0]{leftmargin=12pt,itemindent=0pt,labelwidth=0pt}
\setlist{topsep=.5pt}
\renewcommand{\labelenumi}{\textbf{\arabic{enumi}.}}
\renewcommand{\labelenumii}{\textbf{\alph{enumii}.}}
\thispagestyle{firststyle}
%%%%%%%%%%%%%%%%%%%%%%%%%%%%%%%%%%%%%%%%%%

\newcommand{\xmax}{0}
\newcommand{\xmin}{-\xmax}
\newcommand{\xstep}{0}
\newcommand{\ymax}{0}
\newcommand{\ymin}{-\ymax}
\newcommand{\ystep}{0} 

\textbf{Justify your answers and show your work.}

\emph{Complete the following steps for $f(x)=x^{2}-20$, interval
$[1,9]$, and $n=4$.}
\begin{enumerate}[resume]
\item Sketch the graph of the function on the given interval and illustrate
the \emph{right Riemann sums }using $n$ subintervals. Clearly label
the gridpoints.\\
\renewcommand{\xmax}{10}
\renewcommand{\xmin}{-1}
\renewcommand{\xstep}{0} %must be xmin+n
\renewcommand{\ymax}{70}
\renewcommand{\ymin}{-22}
\renewcommand{\ystep}{4} %must be ymin+n
%%%%%%%
\begin{tikzpicture}
	\begin{axis}[axis x line=bottom, axis y line=left,
		%y=.65cm/10,
		%x=.65cm,
		axis x line=middle, axis y line=middle,     		
		%xtick={\xmin,\xstep,...,\xmax},
		%ytick={\ymin,\ystep,...,\ymax},
		minor x tick num={1},
		minor y tick num={0},
		%grid=both,
		xlabel={$$},
		%axis y label/.style={at={(current axis.above origin)},anchor= east},    
		%y label style={at={(-0.1,0.65)}},
		ylabel={$f(x)$},
		xmin=\xmin.25,xmax=\xmax.25,
		ymin=\ymin.25,ymax=\ymax.25,     		width=8cm,     		height=6cm] 
		\addplot[mark=noner,smooth,red,thick,domain=-1:10]{x^2-20}; 
	\end{axis}
\end{tikzpicture}\\
\setlength\PreviewBorder{2mm}
\pgfplotsset{     
integral segments/.code={\pgfmathsetmacro\integralsegments{#1}},     
integral segments=4,   integral/.style args={#1:#2}{         ybar interval,
domain=#1+((#2-#1)/\integralsegments)/2:#2+((#2-#1)/\integralsegments),         samples=\integralsegments+1,         x filter/.code=\pgfmathparse{\pgfmathresult-((#2-#1)/\integralsegments)}     }}   \begin{tikzpicture}[/pgf/declare function={f=x^2-20;}]     \begin{axis}[         domain=0:10,         samples=100,         axis lines=middle     ]     \addplot [         red,         fill=red!50,         integral=2:6     ] {f};     \addplot [thick,red] {f};     \end{axis}   \end{tikzpicture}
\item Calculate the \emph{right Riemann sums }for the given interval and
number of subintervals, $n$.\vspace{2.5cm}
\end{enumerate}
\emph{Evaluate the following using the provided values.}

\begin{multicols}{2}

\emph{
\begin{eqnarray*}
{\displaystyle \int_{0}^{\nicefrac{\pi}{2}}g(x)}\,dx & = & 1\\
{\displaystyle \int_{\nicefrac{\pi}{2}}^{\pi}g(x)}\,dx & = & \mbox{\ensuremath{\pi}-1}\\
{\displaystyle \int_{\pi}^{\nicefrac{3\pi}{2}}g(x)}\,dx & = & \pi+1\\
{\displaystyle \int_{\nicefrac{3\pi}{2}}^{2\pi}g(x)}\,dx & = & 2\pi-1
\end{eqnarray*}
}

\columnbreak
\begin{enumerate}[resume]
\item ${\displaystyle \int_{0}^{\pi}g(x)}\,dx$\vspace{1cm}\vfill
\item ${\displaystyle \int_{\pi}^{0}\frac{g(x)}{\pi}}\,dx$\vspace{1cm}\vfill
\item Suppose that $g(x)$ is \emph{even}, \\
find ${\displaystyle \int_{-\nicefrac{\pi}{2}}^{\nicefrac{\pi}{2}}g(x)}\,dx$.\vspace{1cm}\vfill
\end{enumerate}
\end{multicols}
\begin{enumerate}[resume]
\item [\textbf{6.}]\setcounter{enumi}{6}Evaluate the expression${\displaystyle \sum_{n=0}^{2}\left(3n+1\right)}$
by writing it as the sum of each of its terms.\vspace{2cm}
\end{enumerate}
\newpage

\emph{Use geometry or symmetry to evaluate the definite integral.
Recall that the area of a triangle is $A=\frac{1}{2}bh$.}
\begin{enumerate}[resume]
\item ${\displaystyle \int_{-2}^{2}\left(2-|x|\right)\,dx}$\vfill
\end{enumerate}
\emph{Use the value of the first integral $I$ to evaluate the two
given integrals. $I={\displaystyle \int_{a}^{b}\left(x-x^{2}\right)\,dx}=-1$.}
\begin{enumerate}[resume]
\item ${\displaystyle \int_{a}^{b}\left(3x^{2}-3x\right)\,dx}$\vspace{2cm}
\end{enumerate}
\emph{Evaluate the following.}
\begin{enumerate}[resume]
\item Let $a$ and $b$ be constants.

\begin{enumerate}[resume]
\item ${\displaystyle {\displaystyle \frac{d}{dx}\int_{a}^{x}3t^{2}\,dt}}$
\vspace{1cm}
\item ${\displaystyle {\displaystyle \frac{d}{dx}\int_{a}^{b}f(t)\,dt}}$\vspace{1cm}
\end{enumerate}
\item $\int_{0}^{2}(x^{2}+1)\,dx$\vfill
\end{enumerate}
\newpage
\begin{enumerate}[resume]
\item $\int_{0}^{\pi}\frac{1}{2}\sin\frac{\theta}{2}\,d\theta$.\vfill
\item ${\displaystyle \int\frac{1-z^{4}}{z^{2}}}\,dz$\vfill
\item ${\displaystyle \int}x^{^{3}}(x^{4}+4)^{15}\,dx$\vfill
\item $\int\frac{1}{(5x-3)^{2}}\,dx$\vfill
\end{enumerate}
\newpage
\begin{enumerate}[resume]
\item ${\displaystyle \int\cos x\sin^{3}x\,dx}$\vfill
\end{enumerate}
\emph{Recall that $\sin x$ is an }\emph{\uline{odd}} \emph{function. }
\begin{enumerate}[resume]
\item ${\displaystyle {\displaystyle \int_{-5}^{5}}\sin^{3}x\,dx}$\vspace{3.5cm}
\end{enumerate}
\emph{Solve the following word problems. }
\begin{enumerate}[resume]
\item An arch is 10 ft high and a 20 ft base. Its shape can be modeled by
$y=10-\frac{x^{2}}{10}$. Find the average height of the arch above
the ground. Note that the maximum height is at $x=0$.\vfill
\item A man is walking in a straight line away from JCSU at the rate of
$2$ mi/hr. If he is $1$ mile from JCSU when he starts walking ($t=0$),
find a function descrinbing his distance from JCSU. \vfill
\end{enumerate}

\end{document}
